\begin{table}[htbp]
  \centering
  \caption{Durch die Riegelstellung geschützte Fläche je Brennstoffmodell, Szenario S3.}
  \label{tab:mtt-fuel}
  \begin{tabular}{lrrrr}
    \toprule
    Brennstoffmodell & \multicolumn{2}{c}{erreichte Fläche (ha)} & \multicolumn{2}{c}{geschützt} \\
    \cmidrule(lr){2-3}\cmidrule(lr){4-5}
     & ohne Barriere & mit Barriere & (ha) & (\%) \\
    \midrule
    Laubwald (TL6)   & 37{,}09 & 22{,}73 & 14{,}42 & 43{,}1 \\
    Heide (GA-heath) & 22{,}95 & 12{,}46 & 10{,}50 & 31{,}4 \\
    Gras (GR4)       & 14{,}64 &  8{,}69 &  5{,}96 & 17{,}8 \\
    Nadelwald (TU5)  &  8{,}43 &  5{,}85 &  2{,}57 &  7{,}7 \\
    \midrule
    Summe            & 83{,}10 & 49{,}73 & 33{,}44 & 100{,}0 \\
    \bottomrule
  \end{tabular}
  \par\smallskip\footnotesize Die Prozentangabe bezeichnet den Anteil des jeweiligen Brennstoffmodells an der insgesamt geschützten Fläche. Bezogen auf den im Untersuchungsgebiet vorhandenen Bestand sinkt der verbrannte Flächenanteil durch die Riegelstellung in der Heide von 95 auf 52\,\%, im Gras von 95 auf 56\,\%, im Nadelwald von 73 auf 51\,\% und im Laubwald von 65 auf 40\,\%.
\end{table}
